\par\indent\SourceLeaf{241}%
Combining this equation with equation (10) gives
\begingroup\fontsize{10}{12}\selectfont
\[
\begin{aligned}
DD''-D'D'&=(\alpha\alpha''+\beta\beta''+\gamma\gamma''-\alpha'\alpha'-\beta'\beta'-\gamma'\gamma')\Delta\\
&\quad+E(n'n'-nn'')+F(nm''-2m'n'+mn'')+G(m'm'-mm'')
\end{aligned}
\]
\endgroup
Now it is evident that
\[
\begin{gathered}
\frac{\mathrm{d}E}{\mathrm{d}p}=2m,\quad\frac{\mathrm{d}E}{\mathrm{d}q}=2m',\quad\frac{\mathrm{d}F}{\mathrm{d}p}=m'+n,\\
\frac{\mathrm{d}F}{\mathrm{d}q}=m''+n',\quad\frac{\mathrm{d}G}{\mathrm{d}p}=2n',\quad\frac{\mathrm{d}G}{\mathrm{d}q}=2n''
\end{gathered}
\]
or
\[
\begin{array}{lll}
m=\tfrac{1}{2}\dfrac{\mathrm{d}E}{\mathrm{d}p},&m'=\tfrac{1}{2}\dfrac{\mathrm{d}E}{\mathrm{d}q},&m''=\dfrac{\mathrm{d}F}{\mathrm{d}q}-\tfrac{1}{2}\dfrac{\mathrm{d}G}{\mathrm{d}p}\\[6pt]
n=\dfrac{\mathrm{d}F}{\mathrm{d}p}-\tfrac{1}{2}\dfrac{\mathrm{d}E}{\mathrm{d}q},&n'=\tfrac{1}{2}\dfrac{\mathrm{d}G}{\mathrm{d}p},&n''=\tfrac{1}{2}\dfrac{\mathrm{d}G}{\mathrm{d}q}
\end{array}
\]
It is also easy to verify that
\[
\begin{aligned}
&\alpha\alpha''+\beta\beta''+\gamma\gamma''-\alpha'\alpha'-\beta'\beta'-\gamma'\gamma'\\
&\quad=\frac{\mathrm{d}n}{\mathrm{d}q}-\frac{\mathrm{d}n'}{\mathrm{d}p}=\frac{\mathrm{d}m''}{\mathrm{d}p}-\frac{\mathrm{d}m'}{\mathrm{d}q}\\
&\quad=-\tfrac{1}{2}\mathord{.}\frac{\mathrm{d}\mathrm{d}E}{\mathrm{d}q^2}+\frac{\mathrm{d}\mathrm{d}F}{\mathrm{d}p\mathord{.}\mathrm{d}q}-\tfrac{1}{2}\mathord{.}\frac{\mathrm{d}\mathrm{d}G}{\mathrm{d}p^2}
\end{aligned}
\]
If we now substitute these various expressions into the formula for the measure of curvature derived at the end of the preceding article, we arrive at the following formula, which is formed solely from the quantities \(E,F,G\) and their first- and second-order differential quotients:
\[
\begin{aligned}
4(EG-FF)^2k&=E(\tfrac{\mathrm{d}E}{\mathrm{d}q}\mathord{.}\tfrac{\mathrm{d}G}{\mathrm{d}q}-2\tfrac{\mathrm{d}F}{\mathrm{d}p}\mathord{.}\tfrac{\mathrm{d}G}{\mathrm{d}q}+(\tfrac{\mathrm{d}G}{\mathrm{d}p})^2)\\
&\quad+F(\tfrac{\mathrm{d}E}{\mathrm{d}p}\mathord{.}\tfrac{\mathrm{d}G}{\mathrm{d}q}-\tfrac{\mathrm{d}E}{\mathrm{d}q}\mathord{.}\tfrac{\mathrm{d}G}{\mathrm{d}p}-2\tfrac{\mathrm{d}E}{\mathrm{d}q}\mathord{.}\tfrac{\mathrm{d}F}{\mathrm{d}q}\\
&\qquad+4\tfrac{\mathrm{d}F}{\mathrm{d}p}\mathord{.}\tfrac{\mathrm{d}F}{\mathrm{d}q}-2\tfrac{\mathrm{d}F}{\mathrm{d}p}\mathord{.}\tfrac{\mathrm{d}G}{\mathrm{d}p})\\
&\quad+G(\tfrac{\mathrm{d}E}{\mathrm{d}p}\mathord{.}\tfrac{\mathrm{d}G}{\mathrm{d}p}-2\mathord{.}\tfrac{\mathrm{d}E}{\mathrm{d}p}\mathord{.}\tfrac{\mathrm{d}F}{\mathrm{d}q}+(\tfrac{\mathrm{d}E}{\mathrm{d}q})^2)\\
&\quad-2(EG-FF)(\tfrac{\mathrm{d}\mathrm{d}E}{\mathrm{d}q^2}-2\tfrac{\mathrm{d}\mathrm{d}F}{\mathrm{d}p\mathord{.}\mathrm{d}q}+\tfrac{\mathrm{d}\mathrm{d}G}{\mathrm{d}p^2})
\end{aligned}
\]
\Article{12}
Since the following holds without restriction
\[
\mathrm{d}x^2+\mathrm{d}y^2+\mathrm{d}z^2=E\mathrm{d}p^2+2F\mathrm{d}p\mathord{.}\mathrm{d}q+G\mathrm{d}q^2
\]
it is evident that \(\BareRoot{(E\mathrm{d}p^2+2F\mathrm{d}p\mathord{.}\mathrm{d}q+G\mathrm{d}q^2)}\) is the general expression for the line element on a curved surface. Thus the analysis developed in the preceding article shows that, to find the measure of curvature, we need no finite formulas expressing the coordinates
